With the rise of microservices, multi-service applications became the de-facto standard for delivering enterprise IT applications \cite{Soldani2018_MicroservicesPainsGains}.
Many big players are already delivering their core business through multi-service applications, with Amazon, Netflix, Spotify, and Twitter perhaps being the most prominent examples.
Multi-service applications, like microservice-based applications, hold the promise of exploiting the potentials of cloud computing to obtain cloud-native applications, \viz composed by loosely coupled services that can be indepedently deployed and scaled \cite{Kratzke2017_CloudNativeApplications}.

At the same time, services and service interactions proliferate in modern multi-service applications, often consisting of hundreds of interacting services.
This makes it harder to monitor the various services forming an application to detect whether they have failed, as well as to understand whether a service failed on its own or in cascade, \viz because some of the services it interacts with failed, causing the service to fail as well.
%Application operators can inspect the logs stored in the distributed environments (\eg virtual machines or containers) used to run the various services, with the aim of identifying whether a service failed, or whether such failure caused/was caused by some other failures in other services \cite{01_Zhou2021_FaultAnalysisDebuggingMicroservices}.
%The resulting process is however time-consuming and error prone, hence making the 
The detection and understanding of failures in modern multi-service applications is actually considered a concrete \enquote{pain} by their operators \cite{Soldani2018_MicroservicesPainsGains}.

Various solutions have been proposed to automate the detection of failures and to automatically determine their possible root causes.
%The idea can be well-explained with an analogy: when a patient feels sick, her doctor looks at her symptoms to diagnose whether she is ill; if this is the case, the doctor prescribes to her patient further analyses to determine the actual root causes of her illness.
%Following the same idea, e
Existing solutions for failure detection rely on identifying anomalies in the behaviour of services, which can be symptoms of their possible failures \cite{38_Nandi2016_AnomalyDetectionControlFlowGraphMining,34_Nedelkoski2019_AnomalyDetectionDistributedTracing,09_Wu2020_MicroserviceDiagnosisDeepLearning}.
Once an anomaly has been detected in a multi-service application, further analyses are enacted to determine the possible root causes for such an anomaly \cite{06_Lin2018_Microscope,49_Shan2019_epsilonDiagnosis}.
This allows application operators to determine whether the anomaly on a service was due to the service itself, to other services underperforming or failing as well, or to environmental reasons, \eg unforeseen peaks in user requests \cite{42_Nguyen2011_PAL} or lack of computing resources in the runtime environment \cite{40_Chen2020_CauseInfer}.

Existing solutions for anomaly detection and root cause analysis are however scattered across different pieces of literature, and often focus only on either anomaly detection or root cause analysis.
This hampers the work of application operators wishing to equip their multi-service applications with a pipeline for detecting anomalies and identifying their root causes.
To this end, in this article we survey the existing techniques for detecting anomalies in modern multi-service applications and for identifying the possible root causes of detected anomalies.
To further support an application operator in choosing the techniques most suited for her application, we also discuss the instrumentation needed to apply an anomaly detection/root cause analysis technique, as well as the additional artifacts that must be provided as input to such techniques.
We also highlight whether the surveyed techniques already put anomaly detection and root cause analysis in a pipeline, or whether they need to be integrated one another.
\upd{In the latter case, we comment on whether the type of anomalies that can be observed with a given anomaly detection technique can be explained by a given root cause analysis technique.}

We believe that our survey can provide benefits to both practitioners and researchers working with modern multi-service applications.
We indeed not only than help them in finding the anomaly detection and root cause analysis techniques most suited to their needs, but we also discuss some open challenges and possible research directions on the topic. 
%The latter set a starting point for researchers wishing to ride the wave of research on modern multi-service applications, by providing them with concrete problems to be solved to further support practitioners daily using them.